%LTeX: language=fr-FR
\documentclass[12pt,a4paper]{article}
\usepackage[utf8]{inputenc}
\usepackage[T1]{fontenc}
\usepackage[french]{babel}
\usepackage[margin=1.7cm]{geometry}
\usepackage{amsmath,amssymb}
\usepackage{graphicx}
\usepackage{enumitem}
\usepackage{tcolorbox}
\usepackage{xcolor}
\usepackage{multicol}
\usepackage{makecell}

\tcbuselibrary{skins}

\definecolor{bandeau}{HTML}{E8EEF7}
\definecolor{titrebleu}{HTML}{1A4FA0}

% Bandeau de titre pour chaque exercice
\newtcbox{\exobox}{
  colback=bandeau, colframe=bandeau, boxrule=0pt,
  arc=0pt, left=4pt, right=4pt, top=3pt, bottom=3pt,
  fontupper=\bfseries, width=\linewidth, halign=left
}

% \exercice{numéro}{titre}{barème}
\newcommand{\exercice}[3]{%
  \vspace{0.15em}
  \noindent\exobox{Exercice #1 --- #2 \hfill (#3 points)}\par
  \vspace{0.15em}
}

\pagestyle{empty}
\setlist{topsep=2pt, itemsep=3pt, parsep=0pt}
\setlength{\parskip}{2pt}
\setlength{\parindent}{0pt}

% Barème détaillé indiqué entre crochets
\newcommand{\pts}[1]{\hfill\textcolor{titrebleu}{\small[#1]}}

\begin{document}

\begin{center}
  {\LARGE\bfseries Corrigé --- Tableaux croisés et probabilités conditionnelles}\\[0.4em]
  {\itshape 1\textsuperscript{ère} ST2S \quad $\bullet$ \quad Contrôle de 55 minutes}
\end{center}

% ================= Exercice 1 =================
\exercice{1}{Automatismes}{6}
\begin{enumerate}
  \item $1-\dfrac{2}{7}=\dfrac{7}{7}-\dfrac{2}{7}=\dfrac{5}{7}$, donc $A=\dfrac{5}{\frac{5}{7}}=5\times\dfrac{7}{5}=\mathbf{7}$. \pts{2 : 1 pour $\frac{5}{7}$, 1 pour le résultat}
  \item $-4x+3=-15 \iff -4x=-15-3 \iff -4x=-18 \iff x=\dfrac{-18}{-4}=\dfrac{9}{2}=\mathbf{4{,}5}$. \pts{2 : 1 pour les étapes, 1 pour le résultat}
  \item L'augmentation est de $70-50=20$~euros, soit $\dfrac{20}{50}=0{,}4=40\,\%$ du prix de départ. Réponse \textbf{C}. \pts{2 : 1 pour le calcul, 1 pour la réponse}
\end{enumerate}

% ================= Exercice 2 =================
\exercice{2}{Construire un tableau croisé}{5}
\begin{enumerate}
  \item On complète les totaux par additions. \pts{3 : 1 pour les données bien placées, 0,5 par total}
  \begin{center}
  \renewcommand{\arraystretch}{1.2}
  \begin{tabular}{|c|c|c|c|}
  \hline
   & Sport en club & Pas de sport en club & Total \\ \hline
  Collège & 30 & 20 & \textbf{50} \\ \hline
  Lycée & 10 & 40 & \textbf{50} \\ \hline
  Total & \textbf{40} & \textbf{60} & 100 \\ \hline
  \end{tabular}
  \end{center}
  \item $\dfrac{60}{100}=0{,}6=\mathbf{60\,\%}$ des élèves ne font pas de sport en club. \pts{1}
  \item Parmi les $40$ élèves qui font du sport en club, $30$ sont collégiens~: \pts{1}

  $\dfrac{30}{40}=\dfrac{3}{4}=0{,}75=\mathbf{75\,\%}$.
\end{enumerate}

% ================= Exercice 3 =================
\exercice{3}{Test de dépistage}{9}
\begin{enumerate}
  \item $75\,\%$ de $40$~: $0{,}75\times 40=30$ patients sont malades avec un test positif. \pts{2 : 1 pour le 30, 1 pour les autres cases}

  Puis $40-30=10$~; $50-30=20$~; $200-40=160$~; $160-20=140$~; $200-50=150$.
  \begin{center}
  \renewcommand{\arraystretch}{1.2}
  \begin{tabular}{|c|c|c|c|}
  \hline
   & Patient malade & Patient non malade & Total \\ \hline
  Test positif & \textbf{30} & \textbf{20} & 50 \\ \hline
  Test négatif & \textbf{10} & \textbf{140} & \textbf{150} \\ \hline
  Total & 40 & \textbf{160} & 200 \\ \hline
  \end{tabular}
  \end{center}
  \item \pts{2 : 0,5 par probabilité}
  \begin{enumerate}[label=\alph*)]
    \item $P(M)=\dfrac{40}{200}=0{,}2=\mathbf{20\,\%}$.
    \item $P(T)=\dfrac{50}{200}=0{,}25=\mathbf{25\,\%}$.
    \item $P_M(T)$ est la probabilité que le test soit positif sachant que le patient est malade. Parmi les $40$ patients malades, $30$ ont un test positif~: $P_M(T)=\dfrac{30}{40}=0{,}75=\mathbf{75\,\%}$.
    \item $M\cap\overline{T}$~: \og le patient est malade \textbf{et} son test est négatif\fg{}. Il y a $10$ patients concernés~: $P(M\cap\overline{T})=\dfrac{10}{200}=0{,}05=\mathbf{5\,\%}$.
  \end{enumerate}
  \item On cherche $P_T(M)$. Parmi les $50$ patients au test positif, $30$ sont malades~: \pts{1}

  $P_T(M)=\dfrac{30}{50}=0{,}6=\mathbf{60\,\%}$.
  \item $M\cup T$~: \og le patient est malade \textbf{ou} son test est positif\fg{}. \pts{2 : 1 pour la phrase, 1 pour le calcul}

  On ne compte qu'une fois les $30$ patients malades au test positif~:

  $P(M\cup T)=\dfrac{40+50-30}{200}=\dfrac{60}{200}=0{,}3=\mathbf{30\,\%}$.
  \item Le test se trompe pour $20$ patients non malades au test positif et $10$ patients malades au test négatif~: \pts{2 : 1 pour le calcul, 1 pour la conclusion}

  $\dfrac{20+10}{200}=\dfrac{30}{200}=0{,}15=15\,\%$. Comme $15\,\%>10\,\%$, l'affirmation est \textbf{fausse}.
\end{enumerate}

\end{document}
